% Source: https://artofproblemsolving.com/community/c3h1380113_try_to_solve_this
% Stable identity: aops-topic-1380113-q01
Let $J(n)$ be the number of the last knight sent when there are $n$ knights. Clearly $J(1)=1$.

\textbf{$n=2k$.} The first trip around sends knights $2,4,\dots,2k$, and the next knight considered is knight $1$, who is skipped. So the process starts again on the $k$ knights $1,3,\dots,2k-1$, beginning by skipping the first of them. The knight in position $j$ of this new list is knight $2j-1$, so $J(2k)=2J(k)-1$.

\textbf{$n=2k+1$.} The first trip sends $2,4,\dots,2k$; then knight $2k+1$ is skipped and knight $1$ is sent. The process continues on the $k$ knights $3,5,\dots,2k+1$, starting by skipping knight $3$. The knight in position $j$ is knight $2j+1$, so $J(2k+1)=2J(k)+1$.

From these, by induction on $n$: if $n=2^m+L$ with $0\le L<2^m$, then $J(n)=2L+1$. (It holds for $n=1$. If $n=2k$ then $k=2^{m-1}+\dfrac{L}{2}$ and $J(n)=2\left(L+1\right)-1=2L+1$; if $n=2k+1$ then $k=2^{m-1}+\dfrac{L-1}{2}$ and $J(n)=2L+1$ as well.)

\begin{itemize}
\item $12=8+4$: the last knight sent is $J(12)=\boxed{9}$.
\item $32=32+0$: $J(32)=\boxed{1}$.
\item $1030=1024+6$: $J(1030)=\boxed{13}$.
\end{itemize}
For example, with $12$ knights the order sent is $2,4,6,8,10,12,3,7,11,5,1,9$.
